\section{Control of the angular modes}
\label{sec:control}

The selections in Lemma~\ref{lem:spectral-surplus} will be mixed by
cycling their basis directions.  We first prove that the angular
operators furnish the required finite-dimensional control.  Throughout
this section, \(m=2s-1\), \(s\geq4\), and \(2\leq l\leq L\).  We equip
each real space \(V_l\) of round spherical harmonics with its round
\(L^2\) inner product and write \(N_l=\dim_{\mathbb R}V_l\).
For an endomorphism \(A\) of \(V_l\), put
\[
 A_{\mathrm{tf}}=A-\frac{\operatorname{tr}A}{N_l}I.
\]
For every orthogonal complex structure \(J\) on \(\mathbb R^{2s}\), the
operator \(D_Ju(x)=du_x(Jx)\) preserves \(V_l\).  It is skew-adjoint,
because its flow consists of round isometries.  In particular,
\((D_J^2)_{\mathrm{tf}}\) is a real symmetric trace-free operator.

\begin{proposition}[Simultaneous control]
\label{prop:control}
Fix positive numbers \(\kappa_l\), \(2\leq l\leq L\).  As \(J\) ranges
over all orthogonal complex structures on \(\mathbb R^{2s}\), the tuples
\begin{equation}
 X_J=\bigl(\kappa_l(D_J^2|_{V_l})_{\mathrm{tf}}\bigr)_{l=2}^L
 \label{eq:control-generators}
\end{equation}
generate, under real linear combinations and brackets, the full Lie
algebra
\[
 \bigoplus_{l=2}^L\mathfrak{sl}(V_l).
\]
\end{proposition}

\begin{proof}
We first prove the assertion in a single degree \(l\).  The nonzero
factor \(\kappa_l\) has no effect on this assertion.  Let
\(\mathfrak a\subset\mathfrak{sl}(V_l)\) be the real Lie algebra
generated by \((D_J^2)_{\mathrm{tf}}\), and set
\[
 W=V_l\otimes_{\mathbb R}\mathbb C,\qquad
 \mathfrak a_{\mathbb C}=\mathfrak a\otimes_{\mathbb R}\mathbb C.
\]
Transpose on \(W\) will always refer to the complex bilinear extension
of the real round inner product.  The single-degree argument first
produces a rank-one symmetric endomorphism of \(W\).  Round conjugation
and brackets then produce all trace-free endomorphisms of \(W\).
Finally, we separate the different degrees to obtain simultaneous
control on their direct sum.

\medskip
\noindent\textbf{Diagonal operators and individual weight blocks.}
Choose orthonormal coordinates
\((x_1,y_1,\ldots,x_s,y_s)\) and write \(z_i=x_i+\mathrm i y_i\).
For \(\theta\in\mathbb R^s\), let \(R_\theta\) rotate coordinate
plane \(i\) by \(\theta_i\), and let \(T_\theta u(x)=u(R_\theta x)\).
The \emph{weight space} of \(\nu\in\mathbb Z^s\) is the character
eigenspace
\[
 W_\nu=\{u\in W:T_\theta u=e^{\mathrm i\sum_i\nu_i\theta_i}u
                    \text{ for every }\theta\}.
\]
Thus \(z_i\) has weight \(e_i\), \(\bar z_i\) has weight \(-e_i\),
and \(D_{J_0}\) acts on \(W_\nu\) by \(\mathrm i\sum_i\nu_i\).
In particular, a monomial of holomorphic degree \(b\) and
antiholomorphic degree \(l-b\) has Hopf charge \(2b-l\), consistently
with Section~\ref{sec:counting}.  Every occurring weight satisfies
\[
 \sum_{i=1}^s|\nu_i|\leq l,\qquad
 l-\sum_{i=1}^s|\nu_i|\in2\mathbb Z.
\]
Every weight on the extreme layer
\[
 \Lambda_l=\left\{\nu\in\mathbb Z^s:
                    \sum_{i=1}^s|\nu_i|=l\right\}
\]
occurs and has a one-dimensional weight space: it is spanned by the
monomial using \(z_i^{\nu_i}\) when \(\nu_i\geq0\) and
\(\bar z_i^{-\nu_i}\) when \(\nu_i<0\).  This monomial is harmonic,
since in each coordinate plane it involves at most one of \(z_i,\bar
z_i\).  It is the only degree-\(l\) monomial of that weight.

More generally, every weight satisfying the displayed bound and parity
condition occurs.  To check its multiplicity, set
\(a=(l-\sum_i|\nu_i|)/2\), a nonnegative integer.
Every monomial of weight \(\nu\) is the corresponding signed monomial
of degree \(\sum_i|\nu_i|\), multiplied by
\(\prod_i(z_i\bar z_i)^{a_i}\) with \(a_i\ge0\) and
\(\sum_i a_i=a\).  The full polynomial weight space therefore has
dimension \(\binom{a+s-1}{s-1}\).
The polynomial Laplacian preserves the weight and maps onto the
corresponding weight space two degrees lower: its adjoint in the
positive monomial inner product is a nonzero scalar multiple of
multiplication by \(\sum_i z_i\bar z_i\), which preserves weight and
is injective.  Thus the harmonic weight multiplicity is
\[
 \binom{a+s-1}{s-1}-\binom{a+s-2}{s-1}
 =\binom{a+s-2}{s-2},
\]
where the second term is zero when \(a=0\).
This standard weight formula also appears in
Lauret--Miatello--Rossetti~\cite[Lemma~3.6]{LauretMiatelloRossetti}.
In particular, interior weight spaces may have dimension greater than
one.  The interpolation below isolates an entire source-to-target
weight block; we use a rank-one conclusion only when both endpoints
are on the extreme layer.

Conjugating by a round rotation sends the family of generators to
itself.  Thus \(\mathfrak a_{\mathbb C}\) is invariant under the round
torus and contains each torus Fourier component of any of its
elements.  Indeed, Fourier projection is an integral of scalar
multiples of torus conjugates, and a finite-dimensional vector
subspace is closed.

For \(\sigma=(\sigma_1,\ldots,\sigma_s)\in\{\pm1\}^s\), let \(J_\sigma\)
rotate coordinate plane \(i\) with sign \(\sigma_i\).  On the
\(\nu\)-weight space, \(D_{J_\sigma}^2\) acts by
\(-(\sum_i\sigma_i\nu_i)^2\).  For \(i\ne j\), the identity
\[
 2^{-s}\sum_{\sigma\in\{\pm1\}^s}
 \sigma_i\sigma_j\left[-\left(\sum_a\sigma_a\nu_a\right)^2\right]
 =-2\nu_i\nu_j
\]
shows that \(\mathfrak a_{\mathbb C}\) contains a diagonal operator
\(A_{ij}\) acting on weight \(\nu\) by \(\nu_i\nu_j+c_{ij}\), where
\(c_{ij}\) is independent of \(\nu\).  The scalar term has no effect
on brackets.

An endomorphism Fourier component of \emph{weight shift} \(\delta\)
satisfies \(T_\theta A T_\theta^{-1}
=e^{\mathrm i\sum_i\delta_i\theta_i}A\); equivalently, it maps each
\(W_\nu\) into \(W_{\nu+\delta}\), with absent weight spaces understood
to be zero.  On its block from \(W_\nu\) to \(W_{\nu+\delta}\), the
operator \(\operatorname{ad}A_{ij}\) acts by
\begin{equation}
 \nu_i\delta_j+\nu_j\delta_i+\delta_i\delta_j.
 \label{eq:control-block-eigenvalues}
\end{equation}
If \(\delta\) has at least three nonzero coordinates, these joint
eigenvalues distinguish all source weights.  In fact, equality for
two sources, with difference \(v\), implies
\[
 v_i\delta_j+v_j\delta_i=0\qquad(i\ne j).
\]
On the support of \(\delta\), the ratios \(v_i/\delta_i\) have
pairwise sums zero; three such coordinates force every ratio to be
zero.  Pairing any coordinate outside the support with one in the
support then gives \(v_i=0\) there as well.

There are only finitely many weights.  Polynomials in the commuting
operators \(\operatorname{ad}A_{ij}\) therefore project onto an
individual source block within a fixed shift component.  More
explicitly, for each competing source choose one coordinate of
\eqref{eq:control-block-eigenvalues} that distinguishes it from the
desired source, and multiply the corresponding linear interpolation
factors.  These operations preserve \(\mathfrak a_{\mathbb C}\).
When source and target belong to \(\Lambda_l\), the isolated nonzero
block is a rank-one directed matrix.

\medskip
\noindent\textbf{A rank-one symmetric operator.}
Regard an extreme weight as \(l\) signed tokens, with no opposite
signs at the same coordinate.  We claim that the algebra contains
the directed matrices for every legal move
\[
 \text{two equal-sign tokens at one coordinate}
 \quad\longleftrightarrow\quad
 \text{one token at each of two other coordinates},
\]
where the two latter signs are arbitrary and both endpoints must
remain in \(\Lambda_l\).  Such a shift has three nonzero coordinates.
It remains to exhibit a generator with a nonzero block for the move.

Here is an explicit simultaneous choice of the needed coefficients.
For three distinct coordinate indices \(i,j,k\), let \(X_a,Y_a\)
denote their real coordinate vectors and set
\[
 JX_i=X_j,\qquad JY_i=X_k,\qquad JY_j=Y_k,
\]
with the three reverse images determined by \(J^2=-I\).
Use the standard complex structure on the remaining coordinate
planes.  This defines an orthogonal complex structure.  For signed
variables \(u_a^\epsilon=x_a+\mathrm i\epsilon y_a\),
\(\epsilon\in\{\pm1\}\), its vector field satisfies
\[
 \begin{split}
 D_Ju_i^\epsilon&=-x_j-\mathrm i\epsilon x_k,\\
 D_Ju_j^\beta&=x_i-\mathrm i\beta y_k,\\
 D_Ju_k^\gamma&=y_i+\mathrm i\gamma y_j.
 \end{split}
\]
For every choice of the three signs, the relevant replacement coefficients are
\[
 \begin{array}{c@{\qquad}c@{\qquad}c@{\qquad}c}
 u_i^\epsilon\longmapsto u_j^\beta&
 u_i^\epsilon\longmapsto u_k^\gamma&
 u_j^\beta\longmapsto u_i^\epsilon&
 u_k^\gamma\longmapsto u_i^\epsilon
 \\[3pt]
 -\frac12&-\frac{\mathrm i\epsilon}{2}&
 \frac12&-\frac{\mathrm i\epsilon}{2}
 \end{array}
\]
For completeness, there is no cancellation with other terms of
\(D_J^2\).  In the independent complex linear coordinates
\(z_a,\bar z_a\), the identity \(D_J^2u=-u\) for linear functions gives,
on homogeneous degree-\(l\) polynomials,
\[
 D_J^2p=-lp+
   \sum_{u,v}(D_Ju)(D_Jv)\,\partial_u\partial_vp.
\]
If
\[
 p=(u_i^\epsilon)^a(u_j^\beta)^b(u_k^\gamma)^c M
\]
is an extreme monomial, where \(M\) uses the other coordinate
planes, the coefficient of the split monomial is
\(\mathrm i\epsilon\,a(a-1)/2\), and that of the merge monomial is
\(-\mathrm i\epsilon\,bc/2\).  They are nonzero whenever the indicated
move is legal.  Indeed the exponent deficits of the target force
the differentiated variables to be precisely the two variables
specified by the move.  Since rotations commute with the Euclidean
Laplacian, \(D_J^2p\) is already harmonic; no harmonic projection
changes these coefficients.  Trace subtraction does not affect this
shift.  Fourier projection and the block isolation above now give
the claimed directed matrices.

We next join \(-le_1\) to \(le_1\) by a simple path of these moves.
For even \(l=2h\), use
\[
 -2he_1\longrightarrow
 -(2h-2)e_1+e_2+e_3\longrightarrow\cdots
 \longrightarrow he_2+he_3
 \longrightarrow\cdots\longrightarrow 2he_1,
\]
splitting negative pairs at coordinate \(1\) in the first half and
merging the tokens at \(2,3\) into positive pairs at \(1\) in the
second half.  This includes \(l=2\), for which the path is
\(-2e_1,e_2+e_3,2e_1\).

For odd \(l=2h+1\), \(h\geq1\), both starting signs admit the path
\[
 \pm(2h+1)e_1\longrightarrow\cdots
 \longrightarrow\pm e_1+he_2+he_3
 \longrightarrow(h-1)e_2+he_3+2e_4.
\]
The last step merges the remaining token at \(1\) and one positive
token at \(2\) into two positive tokens at \(4\).  Reverse the
positive-start path and append it to the negative-start path.
The resulting path is simple: the two branches have opposite
nonzero first coordinates until their common endpoint, and no
branch repeats a weight.  For \(l=3\), that endpoint is \(e_3+2e_4\).
All intermediate weights remain extreme.

Choose monomial bases on these one-dimensional weight spaces and
normalize their directed matrices to matrix units.  Nested brackets
along a simple path give its endpoint unit, by
\([E_{c,b},E_{b,a}]=E_{c,a}\) for distinct vertices.  Hence
\(\mathfrak a_{\mathbb C}\) contains the unit from weight \(-le_1\)
to weight \(le_1\).  If \(v=z_1^l\), torus invariance of the bilinear
round pairing shows that \(v\) pairs only with weight \(-le_1\), and
\[
 v^{\mathsf T}\bar v
   =\int_{S^{2s-1}}|z_1|^{2l}\,d\operatorname{vol}_{g_0}>0.
\]
The endpoint unit is consequently a nonzero multiple of
\(vv^{\mathsf T}\).  Notice also that \(v^{\mathsf T}v=0\).

\medskip
\noindent\textbf{From the rank-one operator to all trace-free matrices.}
All round conjugates of \(vv^{\mathsf T}\) belong to the algebra.
We first verify that the round orbit of \(v\) spans \(W\).
Up to a nonzero scalar, its members are the polynomials
\((a\cdot x)^l\) with
\[
 a\in\mathbb C^{2s}\setminus\{0\},\qquad a\cdot a=0.
\]
To see this, write \(a=b+\mathrm i c\) with \(b,c\) real.  The
isotropy condition says that \(b,c\) are orthogonal and have equal
positive length.  After rescaling, this ordered pair is the image
of the first coordinate pair under a real special orthogonal
transformation.

On homogeneous degree-\(l\) polynomials use the Fischer bilinear
form
\[
 (p,q)_{\mathrm F}=p(\partial)q
       =\sum_{|\alpha|=l}\alpha!\,p_\alpha q_\alpha.
\]
Its restriction to real harmonic polynomials is positive definite,
so its complex bilinear restriction to \(W\) is nondegenerate.
If \(p\in W\) annihilates the orbit powers, then
\[
 (p,(a\cdot x)^l)_{\mathrm F}=l!\,p(a)
\]
shows that \(p\) vanishes on the isotropic quadric.  Therefore
\(p\) is divisible by the quadratic polynomial
\(Q(x)=\sum_{a=1}^s(x_a^2+y_a^2)\).  An elementary verification
is to divide by this monic polynomial in the last variable \(y_s\);
the remainder has the form \(A(x')y_s+B(x')\), where \(x'\) denotes
the remaining variables.  Evaluating at the two
distinct roots for generic \(x'\) forces both \(A\) and \(B\) to
vanish identically.

Such a harmonic multiple must be zero.  Indeed multiplication by
\(Q\) is adjoint to the Euclidean Laplacian in the
positive Hermitian Fischer form.  If
\(p=Qq\) and \(\Delta p=0\), then
\[
 \|p\|_{\mathrm F}^2
   =\langle Qq,p\rangle_{\mathrm F}
   =\langle q,\Delta p\rangle_{\mathrm F}=0.
\]
Thus no nonzero element annihilates the orbit, proving its span.

For two orbit vectors \(u,w\),
\[
 [uu^{\mathsf T},ww^{\mathsf T}]
  =(u^{\mathsf T}w)
        (uw^{\mathsf T}-wu^{\mathsf T}).
\]
For fixed \(u\ne0\), the factor \(u^{\mathsf T}w\) is not
identically zero on the orbit, by its spanning property and
nondegeneracy of the round pairing.  It is a real-analytic function
of the round rotation.  The real special orthogonal group is
connected, so its nonzero locus is dense.  Division by this factor
and then passage to limits show that
\(\mathfrak a_{\mathbb C}\) contains
\(uw^{\mathsf T}-wu^{\mathsf T}\) for every pair of orbit vectors.
Their span is the full complex skew-symmetric algebra, since the
orbit spans \(W\).

For every real skew-symmetric \(K\), the operator
\(\exp(\operatorname{ad}K)\) preserves \(\mathfrak a_{\mathbb C}\);
this uses only bracket closure in the complex algebra.
It follows that
\(\mathfrak a_{\mathbb C}\) is invariant under conjugation by the
full real group \(SO(V_l)\).  Here one may use that every special
orthogonal matrix is a product of plane rotations.  This group on
the harmonic-function space sends \(v\), up to a scalar, to every
nonzero isotropic vector of \(W\): their real and imaginary parts
are orthogonal equal-length pairs, and \(\dim V_l\geq3\).
Consequently the algebra contains \(ww^{\mathsf T}\) for every
isotropic \(w\in W\).

These matrices span all trace-free symmetric matrices.  In a real
orthonormal basis, the rank-one matrices associated with
\(e_i+\mathrm i e_j\) and \(e_i-\mathrm i e_j\) have sum
\(2(E_{ii}-E_{jj})\) and difference
\(2\mathrm i(E_{ij}+E_{ji})\).  Together with the skew-symmetric
matrices already obtained, they span \(\mathfrak{sl}(W)\).
We have proved
\(\mathfrak a_{\mathbb C}=\mathfrak{sl}(W)\).
Since \(\mathfrak a\subset\mathfrak{sl}(V_l)\) is a real vector
subspace and complexification preserves dimension, it follows
that \(\mathfrak a=\mathfrak{sl}(V_l)\).

\medskip
\noindent\textbf{Independence of the degrees.}
Let \(\mathfrak h\) be the Lie algebra generated by the tuples
\eqref{eq:control-generators}.  Every one-factor projection is full
by what we have just proved.  Moreover
\[
 N_l=\binom{m+l}{m}-\binom{m+l-2}{m}
     =\binom{m+l-1}{m-1}+\binom{m+l-2}{m-1}
\]
is strictly increasing in \(l\).  The real Lie algebras
\(\mathfrak{sl}(V_l)\) are therefore simple and pairwise
nonisomorphic.  Simplicity can be seen directly: in any nonzero
ideal, a bracket with a matrix unit first produces an element with
an off-diagonal entry if necessary; polynomials in the adjoint
action of a generic trace-free diagonal matrix isolate an
off-diagonal unit.  Brackets with the other matrix units then
produce all off-diagonal units and all diagonal differences.

We give the resulting Lie-algebra Goursat argument explicitly; for the
two-factor formulation, see \cite[Section~3]{FigueroaUngureanu2016}.
Inductively,
suppose the projection onto a preceding product \(\mathfrak b\)
is full, and let \(\mathfrak c\) be the next simple factor.
The projected algebra
\(\mathfrak h'\subset\mathfrak b\oplus\mathfrak c\) is surjective
onto both sides.  The kernel of its projection onto \(\mathfrak b\)
is an ideal of \(\mathfrak c\), and hence is either zero or
\(\mathfrak c\).  In the latter case the whole product is present.
In the former case \(\mathfrak h'\) is the graph of a surjective
homomorphism \(\mathfrak b\to\mathfrak c\).
The image of each simple summand of \(\mathfrak b\) is an ideal in
\(\mathfrak c\); a nonzero image would make that summand isomorphic
to \(\mathfrak c\).  This contradicts their distinct dimensions.
The induction proves the claimed full direct sum.
\end{proof}

\begin{remark}
The quantifier over all orthogonal complex structures in
Proposition~\ref{prop:control} includes both orientation classes.
This is used in the averaging over all sign vectors.  In
particular, when \(s=4\), restricting to one class would identify
some complementary sign characters and would not justify that
averaging step.  The explicit complex structure used for the token
moves can itself be completed in either class by reversing it on
an unused coordinate plane.
\end{remark}

For general background on control systems on Lie groups, see
Jurdjevic--Sussmann~\cite{JurdjevicSussmann}.  We give the finite-word
argument with its nonsingular differential directly.

We now pass from the Lie algebra to exact finite products.  For the
integers \(M_l\) in \eqref{eq:selection}, choose the ordered real
orthonormal dwell eigenbasis \(e_{l,1},\ldots,e_{l,N_l}\).
If \(M_l\geq1\), define a signed cycle by
\[
 \Pi_l e_{l,i}=
 \begin{cases}
 e_{l,i+1},&1\leq i<M_l,\\
 (-1)^{M_l-1}e_{l,1},&i=M_l,\\
 e_{l,i},&i>M_l,
 \end{cases}
\]
and set \(\Pi_l=I\) when \(M_l=0\).
The sign of the \(M_l\)-cycle and its extra column sign cancel,
so \(\det\Pi_l=1\).  Thus
\[
 \Pi=(\Pi_l)_{l=2}^L\in
 \mathcal G:=\prod_{l=2}^L SL(V_l).
\]

\begin{corollary}[A finite word with invertible differential]
\label{prop:word}
Let \(d=\sum_{l=2}^L(N_l^2-1)\).
There exist finitely many orthogonal complex structures
\(J_1,\ldots,J_R\), an open neighborhood \(U\) of zero in
\(\mathbb R^d\), and smooth real functions \(c_1,\ldots,c_R\) on
\(U\) such that
\[
 \mathcal W(h)=
   \exp(c_R(h)X_{J_R})\cdots\exp(c_1(h)X_{J_1})
\]
satisfies \(\mathcal W(0)=\Pi\), and
\[
 D\mathcal W(0):\mathbb R^d\longrightarrow T_\Pi\mathcal G
\]
is an isomorphism.  The same assertion holds for any prescribed
target in \(\mathcal G\).
\end{corollary}

\begin{proof}
Let \(\mathcal H\) be the subgroup of \(\mathcal G\) consisting of
finite products of \(\exp(tX_J)\) with arbitrary real \(t\).  Define
\[
 E=\operatorname{span}_{\mathbb R}
       \{\operatorname{Ad}_hX_J:h\in\mathcal H\}.
\]
Conjugation by every generator exponential preserves \(E\).
Differentiating shows that \(E\) is stable under bracketing with
each \(X_J\).  Since \(E\) contains the generators, it contains
their generated Lie algebra, which is the whole Lie algebra of
\(\mathcal G\) by Proposition~\ref{prop:control}.
Choose a basis
\(\operatorname{Ad}_{h_a}X_{K_a}\), \(1\leq a\leq d\), from this
spanning set.  The map
\[
 F(t_1,\ldots,t_d)=
 \bigl(h_d\exp(t_dX_{K_d})h_d^{-1}\bigr)\cdots
 \bigl(h_1\exp(t_1X_{K_1})h_1^{-1}\bigr)
\]
has value \(I\) at zero and differential
\[
 DF(0)(t)=\sum_{a=1}^d t_a\operatorname{Ad}_{h_a}X_{K_a},
\]
which is an isomorphism.  Every factor \(h_a\), and its inverse,
is a fixed finite word of generator exponentials.  Thus every
value of \(F\) is in \(\mathcal H\).

The inverse function theorem gives an identity neighborhood
contained in \(\mathcal H\), making \(\mathcal H\) an open subgroup.
The group \(SL_N(\mathbb R)\) is connected: polar decomposition
reduces this to connectedness of \(SO(N)\) and of the positive
definite determinant-one matrices, the latter being joined to the
identity by \(S^t\), \(0\leq t\leq1\).
Hence \(\mathcal G\) is connected, and its open subgroup
\(\mathcal H\) is all of \(\mathcal G\).

Choose a fixed finite word representing the desired target and
multiply it by \(F\).  The resulting word has that target as its
value at zero and still has invertible differential.  Expanding
the fixed conjugating words gives the asserted form, with fixed
complex structures and smooth real time parameters.
\end{proof}

\begin{remark}
The real times in Corollary~\ref{prop:word} may have either sign.
They will be realized as integrals of bounded smooth bumps \(z\)
in weak metric pulses of the form \(q=1+z/T\).  Both signs are
admissible once \(T\) is sufficiently large, because \(q>0\) and
converges uniformly to \(1\).  Likewise, the class of \(J\) may be
changed while \(q=1\), when the angular metric is independent of
\(J\).  Thus neither the signed times nor the two orientation
classes impose a restriction on the later construction.
\end{remark}
